% Generated draft: requires mathematical and source review.
\documentclass[11pt]{article}
\usepackage[margin=25mm]{geometry}
\usepackage{fontspec}
\setmainfont{Latin Modern Roman}
\usepackage{amsmath,amssymb,mathtools,unicode-math}
\setmathfont{Latin Modern Math}
\usepackage{enumitem,needspace,xurl,hyperref}
\setlength{\parindent}{0pt}
\setlength{\parskip}{5pt}
\setlength{\emergencystretch}{3em}
\newcommand{\N}{\mathbb{N}}
\newcommand{\Z}{\mathbb{Z}}
\newcommand{\Q}{\mathbb{Q}}
\newcommand{\R}{\mathbb{R}}
\newcommand{\C}{\mathbb{C}}
\newcommand{\F}{\mathbb{F}}
\newcommand{\A}{\mathbb{A}}
\newcommand{\PP}{\mathbb{P}}
\newcommand{\E}{\mathbb{E}}
\newcommand{\eps}{\varepsilon}
\newcommand{\dd}{\,\mathrm{d}}
\begin{document}
\section*{UnsolvedMath ID: 30005005}
\textbf{Source problem code:} catalogue entry
\subsection{Definitions and mathematical statement}
The supplied record presents a research goal rather than a fully specified conjecture about combinatorial data. Its excerpt refers to Conjecture 1.2 without reproducing that conjecture. The following gives the conventional real planar formulation of the named Gasca--Maeztu conjecture as an explicit interpretation requiring comparison with the complete report.

For an integer \(n\geq 1\), let
\[
\Pi_n=\{p\in\R[x,y]:\deg p\leq n\},\qquad N_n=\dim\Pi_n=\binom{n+2}{2}.
\]
A set \(X\subset\R^2\) of \(N_n\) distinct points is \emph{\(n\)-poised} if the evaluation map
\[
\operatorname{ev}_X:\Pi_n\longrightarrow\R^X,\qquad p\longmapsto (p(A))_{A\in X},
\]
is bijective. For such a set and each \(A\in X\), its fundamental interpolation polynomial \(p_A\in\Pi_n\) is uniquely determined by
\[
p_A(B)=\begin{cases}1,&B=A,\\0,&B\in X\setminus\{A\}.\end{cases}
\]
An \(n\)-poised set is called a \(GC_n\) set if, for every \(A\in X\), there exist \(c_A\in\R\setminus\{0\}\) and nonconstant real affine linear polynomials \(\ell_{A,1},\ldots,\ell_{A,n}\) such that
\[
p_A=c_A\prod_{j=1}^n\ell_{A,j}.
\]
In this formulation, a \emph{maximal line} of \(X\) is a real affine line containing exactly \(n+1\) points of \(X\).

\textbf{Named conjecture, under the preceding interpretation.} For every integer \(n\geq1\) and every \(GC_n\) set \(X\subset\R^2\), there exists a real affine line \(L\) such that
\[
|X\cap L|=n+1.
\]
\textbf{Research goal recorded here.} Understand how combinatorial properties associated with these interpolation configurations can be exploited to prove the Gasca--Maeztu conjecture and to formulate and establish appropriate generalizations to an arbitrary number of variables. The clean statement specifically mentions combinatorial properties of maximal lines. The original excerpt more broadly discusses combinatorial structures arising from products of linear forms and Stanley--Reisner ideals; it does not specify a particular combinatorial criterion sufficient for the conjecture.

The excerpt attributes to Sauer--Xu, reference [33], an implicit result concerning generators of the vanishing ideal
\[
I(X)=\{f\in\R[x_1,\ldots,x_d]:f(A)=0\text{ for every }A\in X\}
\]
of a geometrically categorized configuration, with generators of degree \(n+1\) that are products of linear forms. Its extracted point-count, generator-count, and ambient-dimension notation is corrupted, so no precise higher-dimensional theorem is asserted here. It also attributes to reference [18] a construction replacing linear forms by variables that produces Stanley--Reisner ideals. The construction and its hypotheses are absent. Consequently, neither a specific higher-dimensional conjecture nor a precise implication from such a combinatorial encoding is determined by this record.
\subsection{Short English statement}
The recorded goal is to use combinatorial structures in polynomial interpolation, including maximal lines and Stanley--Reisner ideals, to attack the Gasca--Maeztu conjecture and seek generalizations in more variables. Under the conventional real planar interpretation stated above, the conjecture says that every \(GC_n\) configuration has a line through \(n+1\) of its points. The supplied excerpt omits the conjecture itself and contains corrupted higher-dimensional notation, so this interpretation and any proposed generalization require human mathematical review.
\subsection{Sources}
\begin{itemize}
\item Local UnsolvedMath registry: supplied record with permanent ID \(30005005\), source problem code \texttt{catalogue entry}, title ``Combinatorial Approaches to the Gasca--Maeztu Conjecture,'' in the Oberwolfach Reports Open Problems collection. No individual proposer is supplied. The registry supplies proposed year 2022, status \texttt{open}, and statement status \texttt{corrected\_verified}; its assertion of verification against the official report is supplied metadata, not an independent verification of this draft.
\item Supplied report citation: \emph{Mini-Workshop: Interpolation, Approximation, and Algebra} (2023), \emph{Oberwolfach Reports}, pp. 383--428, DOI \texttt{10.4171/owr/2022/7}. \url{https://doi.org/10.4171/owr/2022/7}. The supplied extraction includes a page marker for p. 385. The citation's year 2023 differs from the supplied proposed year and year tag 2022; this discrepancy is unresolved.
\item Supplied literature source: ``Vardanyan, second proof for \(n=5\).'' \url{https://arxiv.org/abs/2108.12799}.
\item Supplied literature source: ``OWR interpolation report.'' \url{https://doi.org/10.4171/owr/2022/7}.
\item Literature assessment dated exactly \texttt{2026-08-22}: the supplied curation labels the problem open and states that the Gasca--Maeztu conjecture is proved through degree \(n=5\), with no known proof or counterexample in the general planar case. This is supplied dated curation, not an independently verified literature assessment; the record itself requests independent verification of time-sensitive status and cited claims.
\item Historical attribution in the supplied excerpt: cases \(n=1,2\) are described as immediate; Busch [8] is credited with the first proofs for \(n=3,4\); [9] and [25] are cited for additional proofs for \(n=4\); and Hakopian--Jetter--Zimmerman [26] are credited with \(n=5\), described relatively as ``last year.'' Further related references are [7, 10, 11, 22, 23]. Sauer--Xu [33] and [18] are cited for the ideal-theoretic and combinatorial background, respectively. Their complete bibliography entries are not supplied, and the relative date cannot be resolved from this excerpt.
\end{itemize}
\end{document}
